\documentclass[11pt]{article}
\usepackage[margin=2.6cm]{geometry}
\usepackage{amsmath,amssymb,amsthm}
\usepackage[colorlinks=true,linkcolor=blue,urlcolor=blue,citecolor=blue]{hyperref}
\newcommand{\OO}{\widetilde{O}}
\newtheorem{theorem}{Theorem}
\setlength{\parskip}{3pt}

\title{An Optimal-Shape FPRAS for \#Knapsack}
\author{Liran Markin}
\date{August 2026}

\begin{document}
\maketitle

\begin{abstract}
\noindent
The fastest known approximation scheme for counting knapsack solutions
is the $\OO(n^{1.5}\varepsilon^{-2})$ FPRAS of Feng and Jin (SODA
2025). We show that one change -- refining their rounding grid by a
factor of $n/\ell$, which their own sampler could not afford but a
level-enumerating sampler gets for free -- drops the total running
time to
\[
\OO\!\left(n^{1.5} + n\,\varepsilon^{-2}\right),
\]
uniformly: the cost of building their data structure, plus the cost of
merely writing down the $\varepsilon^{-2}$ samples any such estimator
reads. Every $\varepsilon$-dependence beyond the sample size is
gone.\footnote{A Lean~4 development accompanying this memo
machine-checks the components introduced here, including the
statistical heart (the overcount bound); see the repository.}
\end{abstract}

\section{The problem and the result}

\#Knapsack asks: given item weights $W_1,\dots,W_n$ and a capacity
$T$, how many of the $2^n$ subsets have total weight at most $T$? The
problem is \#P-hard, so one settles for a randomized
$(1\pm\varepsilon)$-approximation (an FPRAS). After two decades of
progress (Dyer's $\OO(n^{2.5}\varepsilon^{-2})$~\cite{Dyer03} through
several deterministic and randomized improvements), the record is the
$\OO(n^{1.5}\varepsilon^{-2})$ algorithm of Feng and Jin~\cite{FJ25}.

\begin{theorem}[candidate]
\#Knapsack admits an FPRAS running in time
$\OO(n^{1.5} + n\,\varepsilon^{-2})$.
\end{theorem}

Concretely, at $\varepsilon = 10^{-4}$ and $n = 10^{6}$: Feng--Jin
spend $\sim 10^{17}$ operations; the algorithm here spends
$\sim 10^{14}$ -- the full factor $\sqrt{n}\,$ separates them whenever
$\varepsilon \le n^{-1/2}$. For $\varepsilon \ge n^{-1/4}$ the bound
is flat at $\OO(n^{1.5})$, the construction cost alone. The shape is
optimal for this framework: one cannot avoid building the structure,
and any $(1\pm\varepsilon)$ ratio estimator reads
$\Theta(\varepsilon^{-2})$ samples of $\Theta(n)$-ish description
size.

\section{Where their time goes}

Feng--Jin approximate the count as $|\Omega'| \cdot p$, where
$\Omega'$ is the solution set of a \emph{rounded} instance -- counted
exactly by a merge tree of arrays -- and $p = |\Omega|/|\Omega'|$ is
estimated by sampling from $\Omega'$ uniformly and checking true
feasibility. Their time splits into
\[
\underbrace{\OO(n^{1.5})}_{\text{build the tree}}
\;+\;
\underbrace{N \cdot \OO(\ell\sqrt{n})}_{\text{$N$ samples, each a tree
query}},
\qquad
N = \OO\!\left(\frac{n}{\ell\,\varepsilon^{2}}\right),
\]
where $\ell \in [2, 8n]$ is the popular-class parameter of the
instance. The two $\ell$'s cancel: $N$ times the query cost is
$n^{1.5}\varepsilon^{-2}$ regardless of $\ell$.

Why is $N$ so large? Because their grid step is coupled to
$T/\ell^{2}$, coarse enough that the rounded superset overshoots:
$|\Omega'|/|\Omega| = \OO(n/\ell)$, so $p$ can be as small as
$\ell/n$, and estimating a rare ratio to relative $\varepsilon$ costs
$1/(p\varepsilon^{2})$ samples. The overshoot is pure rounding debt:
every mask's weight is known only up to the accumulated rounding
error, and masks inside that error band around $T$ are counted as
solutions whether or not they are.

\section{The observation}

Nothing in the \emph{correctness} of the pipeline forces the coarse
grid. What forces it is their sampler: a query reads array segments,
so its cost scales with array length, and refining the grid by $F$
multiplies every length -- and the whole sampling phase -- by $F$.
The coupling is an artifact of \emph{how they sample}.

In earlier work on this repository we replaced their query by a pruned
witness-split sampler whose per-draw work enumerates the
\emph{witness levels} present at a node -- a quantity bounded by the
subtree's item count -- and never scans array positions. Its
per-sample cost is $\OO(n)$ \emph{independent of the grid}. That
frees the grid entirely:

\begin{quote}
\emph{Make the grid fine enough that the rounded count is honest, and
pay for it only in array length -- which the sampler never reads.}
\end{quote}

Refine the grid by $F = n/\ell$, so the total round-down error of any
mask stays below one band width $T/\ell$. Two things happen at once:
\begin{enumerate}
\item Rounding down never loses a solution, so
$\Omega \subseteq \Omega'$; and every over-counted mask has true
weight within one band width above $T$. Feng--Jin's own band bound
(their Lemma~3.3 -- a statement about \emph{true} weights, blind to
any rounding) says that band carries only a polylogarithmic fraction
of $|\Omega|$. Hence $p = 1/\mathrm{polylog}$ and
$N = \OO(\varepsilon^{-2})$.
\item Construction grows only linearly in $F$: a merge costs
(length $\times$ $\sqrt{\text{levels}}$), and levels do not change.
Summed over the tree the construction is $\OO(F\ell\sqrt{n})$, which
at $F = n/\ell$ is $\OO(n^{1.5})$ -- the parameter $\ell$ cancels
here too, now in our favor.
\end{enumerate}

\section{The algorithm}

\begin{enumerate}
\item Run the Feng--Jin reduction and class decomposition unchanged,
but set every node's grid step a factor $n/\ell$ finer than theirs.
\item Build the merge tree with their witness-count subroutine
(their Theorem~6.1, which is parametric in array length). Alongside
each node, store static per-(level, residual-class) sorted position
arrays with prefix counts.
\item Draw $N = \OO(\varepsilon^{-2})$ samples with the pruned
witness-split sampler: descend the tree splitting each position among
the witness levels of the two children; every mass and rank query is
a prefix lookup.
\item Check each sampled mask's true weight; output $|\Omega'|$ times
the observed feasible fraction (median-amplified).
\end{enumerate}

\section{Why this is $n^{1.5} + n\varepsilon^{-2}$}

Construction: per merge (length $\times \sqrt{\text{levels}}$) at
boosted length $F\ell/2^{h/2}$ and $n_m/2^{h}$ levels sums, over
depths and classes, to $F\ell\sqrt{n}\cdot\mathrm{polylog}
= n^{1.5}\cdot\mathrm{polylog}$ at $F = n/\ell$. Sampling:
$N = \OO(\varepsilon^{-2})$ samples at $\OO(n)$ each --- level
enumeration, grid-independent --- is $\OO(n\varepsilon^{-2})$. There
is no third term.

\section{Discussion}

\textbf{What is verified.} Beyond the sampler's exactness and cost
ledgers, the statistical heart of this result is machine-checked: the
sandwich $\Omega \subseteq \Omega' \subseteq \Omega \cup
\mathrm{band}$ and the overcount bound
$100\,|\Omega'| \le (101 + 200h)\,|\Omega|$, proved in Lean by
composing with our formalization of Feng--Jin's Lemma~3.3.

\textbf{What is inherited.} Their Theorem~6.1 subroutine (used as
published, at longer arrays -- its statement is parametric in
length), the two-phase outer pipeline, and their Lemma~3.3's
instance-side hypotheses, discharged by their own reduction.

\textbf{What remains.} At $\varepsilon = \Theta(1)$ the bound is
$\OO(n^{1.5})$, tying Feng--Jin: the construction itself. Improving
it is the open problem they pose in \S1.3, tied to the
bounded-monotone $(\max,+)$-convolution frontier -- and it is now the
\emph{only} frontier: sampling, for every $\varepsilon$, is settled
at the output-optimal $n\varepsilon^{-2}$. The price paid is space:
the boosted arrays occupy $\OO(n^{1.5})$ words, the same order as
their construction time but more than their unboosted arrays.

\begin{thebibliography}{9}
\bibitem{FJ25} W.~Feng, C.~Jin.
\emph{A Faster Algorithm for Counting Knapsack Solutions}.
SODA 2025. arXiv:2410.22267.
\bibitem{Dyer03} M.~Dyer.
\emph{Approximate counting by dynamic programming}. STOC 2003.
\bibitem{GMW18} P.~Gawrychowski, L.~Markin, O.~Weimann.
\emph{A Faster FPTAS for \#Knapsack}. ICALP 2018.
\end{thebibliography}

\end{document}
